% =====================================================================
%  Résumé (français) et Abstract (anglais)
% =====================================================================
\cleardoublepage
\chapter*{Résumé}
\addcontentsline{toc}{chapter}{Résumé}
\markboth{Résumé}{Résumé}

Les systèmes de génération augmentée par récupération (\anglais{retrieval-augmented
generation}, RAG) appliqués aux documents visuellement riches ont récemment
opéré un changement de paradigme~: plutôt que d'extraire le texte d'une page
par analyse syntaxique ou reconnaissance optique de caractères, ils rendent la
page en pixels et laissent un modèle vision-langage la lire directement. Ce
détour par l'image préserve la mise en page --- quelle cellule appartient à
quelle ligne de tableau, quelle légende à quelle figure --- que la chaîne
d'extraction textuelle détruit. Mais il déplace le problème sans le résoudre~:
il faut toujours décider comment découper une page en unités récupérables, et
les approches existantes le font \emph{à l'aveugle}, sur une grille de hauteur
fixe indifférente au contenu. Un tableau peut être coupé en deux, un paragraphe
scindé au milieu d'une phrase, et le système de récupération doit alors
reconstruire depuis les pixels une segmentation que le moteur de rendu avait
calculée gratuitement quelques millisecondes plus tôt.

Ce travail part d'une observation simple~: pour une page web, cette
segmentation est déjà disponible. Le moteur de mise en page du navigateur
connaît la boîte englobante exacte de chaque élément du DOM avant qu'un seul
pixel ne soit recadré. Nous construisons \textbf{DOM-Graph RAG}, un pipeline
complet qui exploite cette information~: il tuile la page rendue le long des
frontières de ses propres éléments, organise les tuiles obtenues en un graphe
de preuves multigranulaire s'étendant sur plusieurs pages via les liens
hypertextes, entraîne par contraste un modèle vision-langage à évaluer la
pertinence d'une tuile pour une question, et parcourt le graphe sous budget au
moment de la requête pour ne lire que les tuiles utiles.

Le système est implémenté de bout en bout et évalué composant par composant sur
des pages web réelles. Deux résultats sont sans ambiguïté. D'abord, le tuilage
guidé par le DOM améliore radicalement la \emph{trouvabilité} du contenu~: la
proportion de questions dont la réponse survit dans une unité récupérable
passe de \SI{26}{\percent} avec une grille aveugle à \SI{52}{\percent} avec nos
tuiles, et ce à un coût de représentation \emph{inférieur} --- \num{2.7} fois
moins de pixels par page, parce que recadrer un paragraphe sur la largeur
réellement occupée par son texte économise davantage que multiplier les tuiles
ne coûte. Ensuite, l'entraînement contrastif du lecteur fonctionne nettement~:
sur un partitionnement par article strictement disjoint, le Recall@1 passe de
\num{0.433} à \num{0.785} et le MRR de \num{0.578} à \num{0.856}.

Les résultats des étapes en aval sont plus nuancés et nous les analysons en
détail. L'avantage de classement du tuilage DOM sur une grille aveugle est réel
mais modeste une fois corrigé le biais de taille de vivier, et il ne se
maintient pas face à un découpage textuel plat sur le sous-ensemble apparié~;
nous montrons que ce résultat s'explique par le scoreur lexical employé, qui
n'exploite aucune information spatiale. Le contrôleur de preuves budgété
n'établit pas encore de supériorité sur une sélection top-$k$ statique~: nous
diagnostiquons une cause architecturale précise --- la frontière d'exploration
ne franchissait jamais les frontières de document --- implémentons le correctif,
et mesurons une progression de \SI{29.7}{\percent} à \SI{39.6}{\percent} de
couverture inter-pages, contre \SI{40.7}{\percent} pour la référence. Enfin,
l'évaluation bout-en-bout sur des questions multi-sauts humaines révèle que le
mode de défaillance dominant du système assemblé est l'abstention explicite
plutôt que la fabulation, ce qui identifie sans ambiguïté le sous-rappel de
preuves --- et donc le scoreur de pertinence --- comme le verrou restant.

\vspace{0.6em}
\noindent\textbf{Mots-clés} --- génération augmentée par récupération, RAG
visuel, modèles vision-langage, segmentation de pages web, DOM, graphes de
preuves, apprentissage contrastif, LoRA, questions-réponses multi-sauts.

% ---------------------------------------------------------------------
\cleardoublepage
\chapter*{Abstract}
\addcontentsline{toc}{chapter}{Abstract}
\markboth{Abstract}{Abstract}

Retrieval-augmented generation (RAG) over visually rich documents has recently
shifted paradigm: instead of extracting a page's text through parsing or
optical character recognition, systems render the page as pixels and let a
vision-language model read it directly. This detour through the image preserves
the layout information --- which cell belongs to which table row, which caption
to which figure --- that the text-extraction pipeline destroys. But it relocates
the problem rather than solving it: one still has to decide how to cut a page
into retrievable units, and existing approaches do so \emph{blindly}, on a
fixed-height pixel grid indifferent to content. A table can be cut in half, a
paragraph split mid-sentence, and the retriever must then reconstruct from
pixels a segmentation the rendering engine had computed for free milliseconds
earlier.

This work starts from a simple observation: for a web page, that segmentation
already exists. The browser's layout engine knows every DOM element's exact
bounding box before a single pixel is cropped. We build \textbf{DOM-Graph RAG},
a complete pipeline that exploits this: it tiles the rendered page along its own
element boundaries, organizes the resulting tiles into a multigranular evidence
graph spanning several pages via hyperlinks, contrastively trains a
vision-language model to score a tile's relevance to a question, and walks the
graph under a budget at query time so that only useful tiles are read.

The system is implemented end to end and evaluated component by component on
live web pages. Two results are unambiguous. First, DOM-guided tiling
dramatically improves content \emph{findability}: the share of questions whose
answer survives into a retrievable unit rises from \SI{26}{\percent} under a
blind grid to \SI{52}{\percent} with our tiles, and it does so at a
\emph{lower} representation cost --- \num{2.7}$\times$ fewer pixels per page,
because cropping a paragraph to the width its text actually occupies saves more
than multiplying tiles costs. Second, contrastive fine-tuning of the reader
clearly works: on a strictly disjoint article-level split, Recall@1 rises from
\num{0.433} to \num{0.785} and MRR from \num{0.578} to \num{0.856}.

Downstream results are more nuanced and we analyse them in detail. DOM tiling's
ranking advantage over a blind grid is real but modest once pool-size bias is
corrected, and it does not hold against flat text chunking on the paired subset;
we show this follows from the lexical scorer used, which exploits no spatial
information. The budgeted evidence controller does not yet establish
superiority over static top-$k$ selection: we diagnose a precise architectural
cause --- the exploration frontier never crossed document boundaries ---
implement the fix, and measure cross-page coverage rising from \SI{29.7}{\percent}
to \SI{39.6}{\percent} against \SI{40.7}{\percent} for the baseline. Finally,
end-to-end evaluation on human-written multi-hop questions reveals that the
assembled system's dominant failure mode is explicit abstention rather than
confabulation, which unambiguously identifies evidence under-recall --- and
therefore the relevance scorer --- as the remaining bottleneck.

\vspace{0.6em}
\noindent\textbf{Keywords} --- retrieval-augmented generation, visual RAG,
vision-language models, web page segmentation, DOM, evidence graphs,
contrastive learning, LoRA, multi-hop question answering.
